\section[Stage 201]{Stage 201: Realization compiler and canonical orbit projection}
\label{stage:201}

\stagefield{Part / anchor}{Part VI; primary anchor \resultanchor{MTDC-T10}, local subanchor \resultanchor{MTDC-T10.1}.}
\claimstatus{\StatusExactClosure{}.}

\stagefield{Purpose}{Stage~201 turns the reference-free four-scalar home-stretch theorem into an executable realization audit.  It defines intrinsic target ratios, the quotient coordinates \((q_{\rm tr},q_{\rm nt},q_\eta)\), the canonical dependent-triple repair vector, and the same-free-quintuple orbit projection.}

\stagefield{Inputs}{Imports the Stage~192 quotient projector, the Stage~198 finite orbit law, and the Stage~200 packet \((\chi_Q-1,q_{\rm tr},q_{\rm nt},q_\eta)\).  The positive state is \(\mathbf x=(\lambda_W,c_{\eta U},\gamma,K_U,K_\eta^{(\rm eff)},K_W^{(\rm eff)},\mu_W,T_U)\).}

\stagefield{Verification}{SymPy audit: \StageFile{scripts/moving_throat_pde_stage201_explicit_realization_compiler_and_canonical_orbit_projection_sympy_audit.py}.  Mathematica audit: none yet.}

\stagefield{Derivation ledger}{The derivation forms the target monomial ratios, logs them to obtain the quotient coordinates, solves the dependent-triple repair equation \(M_*\Delta\mathbf x_{\rm rep}=-\mathbf q\), exponentiates the repair, and proves the projected state is the unique target-orbit point with the same free quintuple.}

\stagefield{Output}{Intrinsic realization packet, canonical orbit projection \(\Pi_{\mathcal O_*}^{\rm can}\), fixed-point criterion \(\mathbf x\in\mathcal Z_*\iff \chi_Q=1\) and \(\Pi_{\mathcal O_*}^{\rm can}(\mathbf x)=\mathbf x\).}

\stagefield{Verification note}{The detailed theorem-block formulas supporting this card are collected in the Part VI appendix narrative above and in the source stage.  The card restores original stage order and records the claim boundary; it should not be read as a second independent proof.}

\stagefield{Checks}{\begin{verificationchecklist}
\item Check target monomial ratios are formed against the declared target values before taking logarithms.
\item Check the same-free-quintuple convention is used; changing the free/dependent split changes the projection chart.
\item Check graph-aligned closure is not confused with actual PDE realization.
\item Carry the \StatusExactClosure{} status tag whenever citing the compiler itself.
\end{verificationchecklist}}

\stagefield{Downstream use}{This stage feeds the Part VI realization compiler and finite mixed-ray search.  Downstream papers should cite \resultanchor{MTDC-T10.1} for this local stage range and \resultanchor{MTDC-T10} for the full Part VI compiler.}
